% SOURCE_AUTHORITY: sga5_fr_workpass.tex, lines 11799--11843 (LF-normalized unit).
% SOURCE_SHA256: 791F4EFFC5E02832D5D77ED03518C8156D6F07E4C8238B03545DB93D883FBB28
\subsection*{4. Comparación con el caso de las categorías abelianas}

Sea $\cC$ una categoría abeliana y $D(\cC)$ la categoría derivada de $\cC$; se denota, como de costumbre, por $D^b(\cC)$ la subcategoría plena de $D(\cC)$ generada por los complejos acotados. Asociando a cada objeto $X$ de $\cC$ el complejo cuyas componentes son todas nulas, salvo la de grado cero, que es igual a $X$, se obtiene un homomorfismo
\[
K(\cC)\xrightarrow{\mu}K(D^b(\cC)).
\]
Este homomorfismo es, de hecho, un isomorfismo. En efecto, podemos utilizar la proposición 3.1 para construir un inverso de $\mu$, teniendo en cuenta que
\[
D^b(\cC)\simeq K^{b,b}(\cC)/K^{b,\phi}(\cC),
\]
siendo $K^{b,\phi}(\cC)$ la subcategoría gruesa de $K^{b,b}(\cC)$ de los complejos acíclicos. Sea $X^\bullet=(X_i)\in\Ob(K^{b,b}(\cC))$. Si se pone
\[
f(X^\bullet)=\sum_i(-1)^i\operatorname{cl}_{\cC}(X_i),
\]
se obtiene una función aditiva de $K^{b,b}(\cC)$ con valores en $K(\cC)$. En efecto, se observa primero que, si $X^\bullet$ e $Y^\bullet$ son homotópicos, entonces $f(X^\bullet)=f(Y^\bullet)$, puesto que $f(X^\bullet)=f(H(X^\bullet))$. Resulta, por tanto, un homomorfismo
\[
K(D^b(\cC))\xrightarrow{\nu}K(\cC),
\]
y es inmediato comprobar que es el inverso de $\mu$.

Como el invariante $K(\cC)$ carece a menudo de interés cuando $\cC$ es «demasiado grande», conviene relativizarlo respecto de una subcategoría $\cC_o$ de $\cC$. Sea, pues, $\cC_o$ una subcategoría plena de $\cC$. Podemos definir, por una parte, el grupo
\[
K_o(\cC)=K_{\cC_o}(\cC).
\]
Supongamos, por otra parte, que se cumple la condición siguiente:
\[
(A)\qquad \cC_o\text{ es estable bajo sumas directas finitas.}
\]
Es entonces fácil definir la categoría triangulada $D_o^b(\cC)$ utilizando los complejos (de cohomología acotada) de $\cC$ cuyas componentes son objetos de la subcategoría $\cC_o$, y pasando al cociente por los cuasi-isomorfismos de tales complejos. Resulta, como antes, un homomorfismo
\[
(*)\qquad K_o(\cC)\longrightarrow K(D_o^b(\cC)).
\]
Supongamos que la subcategoría $\cC_o$ satisface además la condición (B) siguiente:
\[
(B)\qquad
\text{Si }X,Y\in\Ob\cC_o\text{ y }X\xrightarrow{u}Y
\text{ es un epimorfismo (en }\cC\text{, naturalmente)},
\]
entonces $\operatorname{Ker}u\in\Ob\cC_o$. Razonando como antes, puede demostrarse que el homomorfismo $(*)$ es un isomorfismo.

El funtor canónico de inclusión
\[
D_o^b(\cC)\longrightarrow D^b(\cC)
\]
no es, en general, plenamente fiel. Sin embargo:
